% Part: first-order-logic
% Chapter: natural-deduction
% Section: proving-things-quant

\documentclass[../../../include/open-logic-section]{subfiles}

\begin{document}

\olfileid{fol}{ntd}{prq}

\olsection{\usetoken{P}{derivation} مع المكممات}

\begin{ex}
لا بد مع المكممات من حفظ شرط المتغير المميَّز؛ وقد يدعو حفظه إلى تغيير ترتيب الاستدلالات.
والأنفع في الغالب أن نبدأ بالقواعد المقيدة به عند البحث، فتكون أدنى موضعًا في البرهان التام.

لنر كيف نعطي !!a{derivation} لل!!{formula}
$\lexists[x][\lnot !A(x)] \lif \lnot \lforall[x][!A(x)]$.
نبدأ كالمعتاد، فنكتب:
\begin{prooftree}
\AxiomC{}
\UnaryInfC{$\lexists[x][\lnot !A(x)]\lif \lnot \lforall[x][!A(x)]$}
\end{prooftree}
نبدأ بكتابة ما يلزم لتسويغ الخطوة الأخيرة باستعمال قاعدة~\Intro{\lif}.
\begin{prooftree}
\AxiomC{$\Discharge{\lexists[x][\lnot !A(x)]}{1}$}
\DeduceC{$\lnot \lforall[x][!A(x)]$}
\DischargeRule{\Intro{\lif}}{1}
\UnaryInfC{$\lexists[x][\lnot !A(x)]\lif \lnot \lforall[x][!A(x)]$}
\end{prooftree}
لا تظهر قاعدة نعملها مباشرة في~$\lnot \lforall[x][!A(x)]$؛ فنهيئ ال!!{derivation} لاستعمال~\Elim{\lexists}.
وهنا يلزم اختيار ثابت لا يقع في~$\lexists[x][\lnot !A(x)]$ ولا في فرض يعتمد عليه الاستدلال، صيانةً لشرط المتغير المميَّز.
وحيث لا تظهر !!{constant}s أصلًا، فكل اختيار صالح.
\begin{prooftree}
\AxiomC{$\Discharge{\lexists[x][\lnot !A(x)]}{1}$}
\AxiomC{$\Discharge{\lnot !A(a)}{2}$}
\DeduceC{$\lnot \lforall[x][!A(x)]$}
\DischargeRule{\Elim{\lexists}}{2}
\BinaryInfC{$\lnot \lforall[x][!A(x)]$}
\DischargeRule{\Intro{\lif}}{1}
\UnaryInfC{$\lexists[x][\lnot !A(x)] \lif \lnot \lforall[x][!A(x)]$}
\end{prooftree}
ولإثبات~$\lnot \lforall[x][!A(x)]$ نجرب~\Intro{\lnot}.
فنطلب تناقضًا، مع جواز أخذ
$\lforall[x][!A(x)]$ فرضًا زائدًا. ويجوز أن يستعمل التناقض الفرض~$\lnot !A(a)$، إذ سيسقطه~\Elim{\lexists}.
وترتيب ذلك هكذا:
\begin{prooftree}
\AxiomC{$\Discharge{\lexists[x][\lnot !A(x)]}{1}$}
\AxiomC{$\Discharge{\lnot !A(a)}{2}, \Discharge{\lforall[x][!A(x)]}{3}$}
\DeduceC{$\lfalse$}
\DischargeRule{\Intro{\lnot}}{3}
\UnaryInfC{$\lnot \lforall[x][!A(x)]$}
\DischargeRule{\Elim{\lexists}}{2}
\BinaryInfC{$\lnot \lforall[x][!A(x)]$}
\DischargeRule{\Intro{\lif}}{1}
\UnaryInfC{$\lexists[x][\lnot !A(x)]\lif \lnot \lforall[x][!A(x)]$}
\end{prooftree}
وقد قرب التناقض. أسهل ما نعمله الآن~\Elim{\lforall}، ولا شرط فيه للمتغير المميَّز.
يجوز إحلال أي حد مناسب محل~$x$ المكمم كليًا، فنختار~$a$ ليتحقق التناقض:
\begin{prooftree}
\AxiomC{$\Discharge{\lexists[x][\lnot !A(x)]}{1}$}
\AxiomC{$\Discharge{\lnot !A(a)}{2}$}
\AxiomC{$\Discharge{\lforall[x][!A(x)]}{3}$}
\RightLabel{\Elim{\lforall}}
\UnaryInfC{$!A(a)$}
\RightLabel{\Elim{\lnot}}
\BinaryInfC{$\lfalse$}
\DischargeRule{\Intro{\lnot}}{3}
\UnaryInfC{$\lnot \lforall[x][!A(x)]$}
\DischargeRule{\Elim{\lexists}}{2}
\BinaryInfC{$\lnot \lforall[x][!A(x)]$}
\DischargeRule{\Intro{\lif}}{1}
\UnaryInfC{$\lexists[x][\lnot !A(x)]\lif \lnot \lforall[x][!A(x)]$}
\end{prooftree}

والآن نراجع شرط المتغير المميَّز، كما ينبغي في براهين المكممات.
لم نستعمل من القواعد المقيدة به إلا~\Elim{\exists}، والمتغير المميَّز~$a$ لا يقع في الفروض التي يعتمد عليها الاستدلال.
فهذا !!{derivation} صحيح.
\end{ex}

\begin{ex}
وقد يكون المطلوب اشتقاق !!a{formula} من !!{formula}s أخرى، فتبقى عندئذ فروض غير مسقطة.
فلا يكفي ضبط النتيجة المرادة؛ بل لا بد من ضبط الفروض معها.

لنر كيف نعطي !!a{derivation} لل!!{formula}~$\lexists[x][!C(x,b)]$
من الافتراضين~$\lexists[x][(!A(x) \land !B(x))]$ و
$\lforall[x][(!B(x) \lif !C(x,b))]$. نبدأ كالمعتاد بكتابة النتيجة في
الأسفل.
\begin{prooftree}
\AxiomC{}
\UnaryInfC{$\lexists[x][!C(x,b)]$}
\end{prooftree}

لنا مقدمتان. فإن استعملنا الأولى لطلب
!!a{derivation} لـ~$\lexists[x][!C(x, b)]$ من
$\lexists[x][(!A(x) \land !B(x))]$، كان الطريق~\Elim{\lexists}.
ولأن لهذه القاعدة شرط المتغير المميَّز، نقدمها في البحث، فنحصل على:
\begin{prooftree}
\AxiomC{$\lexists[x][(!A(x) \land !B(x))]$}
\AxiomC{$\Discharge{!A(a) \land !B(a)}{1}$}
\DeduceC{$\lexists[x][!C(x,b)]$}
\DischargeRule{\Elim{\lexists}}{1}
\BinaryInfC{$\lexists[x][!C(x,b)]$}
\end{prooftree}
يشترك الافتراضان اللذان نعمل بهما في~$!B$. وقد يكون من المفيد في هذه
المرحلة أن نطبق~\Elim{\land} لفصل~$!B(a)$.
\begin{prooftree}
\AxiomC{$\lexists[x][(!A(x) \land !B(x)])$}
\AxiomC{$\Discharge{!A(a) 
\land !B(a)}{1}$}
\RightLabel{\Elim{\land}}
\UnaryInfC{$!B(a)$}
\DeduceC{$\lexists[x][!C(x,b)]$}
\DischargeRule{\Elim{\lexists}}{1}
\BinaryInfC{$\lexists[x][!C(x,b)]$}
\end{prooftree}

وأما الفرض الثاني فهو~$\lforall[x][(!B(x) \lif
!C(x,b))]$. لا شرط هنا للمتغير المميَّز؛ فنضع محل~$x$ ال!!{constant}~$a$ بقاعدة~\Elim{\lforall}، فينتج
$!B(a) \lif !C(a, b)$.
فقد اجتمع~$!B(a) \lif !C(a,b)$ و$!B(a)$، والخطوة التالية تطبيق~\Elim{\lif} مباشرة:
\begin{prooftree}
\AxiomC{$\lexists[x][(!A(x) \land !B(x))]$}
\AxiomC{$\lforall[x][(!B(x) \lif !C(x,b))]$}
\RightLabel{\Elim{\lforall}}
\UnaryInfC{$!B(a) \lif !C(a,b)$}
\AxiomC{$\Discharge{!A(a) 
\land !B(a)}{1}$}
\RightLabel{\Elim{\land}}
\UnaryInfC{$!B(a)$}
\RightLabel{\Elim{\lif}}
\BinaryInfC{$!C(a,b)$}
\DeduceC{$\lexists[x][!C(x,b)]$}
\DischargeRule{\Elim{\lexists}}{1}
\insertBetweenHyps{\hspace{-5em}}
\BinaryInfC{$\lexists[x][!C(x,b)]$}
\end{prooftree}
لقد أوشكنا!{} فبتطبيق واحد لـ~\Intro{\lexists} نبلغ غايتنا.
\begin{prooftree}
\AxiomC{$\lexists[x][(!A(x) \land !B(x))]$}
\AxiomC{$\lforall[x][(!B(x) \lif !C(x,b))]$}
\RightLabel{\Elim{\lforall}}
\UnaryInfC{$!B(a) \lif !C(a,b)$}
\AxiomC{$\Discharge{!A(a) 
\land !B(a)}{1}$}
\RightLabel{\Elim{\land}}
\UnaryInfC{$!B(a)$}
\RightLabel{\Elim{\lif}}
\BinaryInfC{$!C(a,b)$}
\RightLabel{\Intro{\lexists}}
\UnaryInfC{$\lexists[x][!C(x,b)]$}
\DischargeRule{\Elim{\lexists}}{1}
\insertBetweenHyps{\hspace{-5em}}
\BinaryInfC{$\lexists[x][!C(x,b)]$}
\end{prooftree}
وبما أننا حرصنا في كل خطوة على ألا تُنتهك شروط المتغير المميَّز،
يمكننا الوثوق بأن هذا !!{derivation} صحيح.
\end{ex}

\begin{ex}
أعط !!a{derivation} لل!!{formula}~$\lnot\lforall[x][!A(x)]$ من
الافتراضين~$\lforall[x][!A(x)] \lif \lexists[y][!B(y)]$ و
$\lnot\lexists[y][!B(y)]$. نبدأ كالمعتاد بكتابة ال!!{formula}
المطلوبة في الأسفل.
\begin{prooftree}
\AxiomC{}
\UnaryInfC{$\lnot\lforall[x][!A(x)]$}
\end{prooftree}
السطر الأخير في ال!!{derivation} نفي، ولذلك فلنجرب استعمال
\Intro{\lnot}. وسيتطلب هذا أن نعرف كيف !!{derive} تناقضًا.
\begin{prooftree}
\AxiomC{$\Discharge{\lforall[x][!A(x)]}{1}$}
\DeduceC{$\lfalse$}
\DischargeRule{\Intro{\lnot}}{1}
\UnaryInfC{$\lnot\lforall[x][!A(x)]$}
\end{prooftree}
قد تصلح~\Elim{\lforall}، لكن نفعها للمطلوب غير بيّن.
فلنرجع إلى الفروض: إن~$\lforall[x][!A(x)] \lif \lexists[y][!B(y)]$
ومعه~$\lforall[x][!A(x)]$ يهيئان تطبيق~\Elim{\lif}:
\begin{prooftree}
\AxiomC{$\lforall[x][!A(x)] \lif \lexists[y][!B(y)]$}
\AxiomC{$\Discharge{\lforall[x][!A(x)]}{1}$}
\RightLabel{\Elim{\lif}}
\BinaryInfC{$\lexists[y][!B(y)]$}
\DeduceC{$\lfalse$}
\DischargeRule{\Intro{\lnot}}{1}
\UnaryInfC{$\lnot\lforall[x][!A(x)]$}
\end{prooftree}
وبقي الفرض الأخير؛ وبه نبلغ التناقض باستعمال~\Elim{\lnot}:
\begin{prooftree}
\AxiomC{$\lnot\lexists[y][!B(y)]$}
\AxiomC{$\lforall[x][!A(x)] \lif \lexists[y][!B(y)]$}
\AxiomC{$\Discharge{\lforall[x][!A(x)]}{1}$}
\RightLabel{\Elim{\lif}}
\BinaryInfC{$\lexists[y][!B(y)]$}
\RightLabel{\Elim{\lnot}}
\BinaryInfC{$\lfalse$}
\DischargeRule{\Intro{\lnot}}{1}
\UnaryInfC{$\lnot\lforall[x][!A(x)]$}
\end{prooftree}
\end{ex}

\begin{prob}
أعط !!{derivation}s تبين ما يأتي:
\begin{enumerate}
\item $\Proves (\lforall[x][!A(x)] \land \lforall[y][!B(y)]) \lif
\lforall[z][(!A(z) \land !B(z))]$.
\item $\Proves (\lexists[x][!A(x)] \lor \lexists[y][!B(y)]) \lif
\lexists[z][(!A(z) \lor !B(z))]$.
\item $\lforall[x][(!A(x) \lif !B)] \Proves \lexists[y][!A(y)] \lif !B$.
\item $\lforall[x][\lnot !A(x)] \Proves \lnot\lexists[x][!A(x)]$.
\item $\Proves \lnot\lexists[x][!A(x)] \lif \lforall[x][\lnot !A(x)]$.
\item $\Proves \lnot\lexists[x][\lforall[y][((!A(x,y) \lif \lnot
!A(y,y)) \land (\lnot !A(y,y) \lif !A(x,y)))]]$.
\end{enumerate}
\end{prob}

\begin{prob}
أعط !!{derivation}s تبين ما يأتي:
\begin{enumerate}
\item $\Proves \lnot\lforall[x][!A(x)] \lif \lexists[x][\lnot!A(x)]$.
\item $(\lforall[x][!A(x)] \lif !B) \Proves \lexists[y][(!A(y) \lif !B)]$.
\item $\Proves \lexists[x][(!A(x) \lif \lforall[y][!A(y)])]$.
\end{enumerate}
(تتطلب هذه كلها قاعدة~$\FalseCl$.)
\end{prob}

\end{document}
